\section{Modelo y producto: del modelo desnudo a Cloud/OS/contenedores}

La sección anterior trató el Online Learning; esta sección pasa al negocio y al producto. Guangmi (广密) considera que la era de los lanzamientos de modelos desnudos podría ir perdiendo peso y que las barreras competitivas vendrán cada vez más de Cloud, el OS, el ecosistema y los contenedores de flujo de trabajo. La razón es que otros pueden alcanzar las capacidades de un modelo, el código abierto las presiona y la API las convierte en mercancía; en cambio, el entorno donde el usuario completa sus tareas cada día y el sistema que posee el estado, los datos, los permisos y la relación de pago son lo que forma una barrera duradera.

\begin{figure}[H]
\centering
\includegraphics[width=0.92\textwidth]{model-product-moat.png}
\caption{Barreras de modelo y producto: la era de los lanzamientos de modelos desnudos pierde peso y Cloud/OS/el ecosistema se vuelven la línea de defensa. Diagrama conceptual propio, elaborado a partir de la conversación de 01:02:45--01:15:11.}
\end{figure}

\begin{knowledgebox}{Lectura de la figura: un modelo fuerte no equivale a un producto fuerte}
A la izquierda, Model destaca el techo de capacidad y la eficiencia de costos; a la derecha, Product/OS destaca el punto de entrada, el flujo de trabajo, el ecosistema, la facturación y el retorno de datos. Al leer esta figura hay que retener una idea: el modelo es el motor; el producto es la carretera, el tablero, el sistema de cobro y la experiencia de manejo.
\end{knowledgebox}

\subsection{El precio de 20 dólares y el problema de copiar el SaaS}

Esta subsección examina la fijación de precios. La entrevista pregunta por qué muchos productos de IA cuestan cerca de 20 dólares y si eso solo copia al SaaS. La cuestión implícita de Guangmi es que el costo marginal de un producto de IA es distinto del de un SaaS tradicional. Una tarea compleja de Agent puede consumir muchos token y llamadas a herramientas; si se vende como suscripción mensual barata de uso ilimitado, aparece presión sobre el margen bruto. Por eso, en el futuro, los precios podrían acercarse a un modelo mixto que combine uso, valor de la tarea, licencias empresariales por usuario y consumo de cómputo.

\begin{warningbox}{El precio no es un problema de UI, sino de economía del cómputo}
Si un producto empaqueta tareas de Agent de alto costo como uso ilimitado a bajo precio, puede crecer a corto plazo, pero a largo plazo el costo de inferencia y el abuso de los usuarios lo aplastarán. Un gerente de producto de IA debe entender los token, la caché, las llamadas a herramientas y el valor de la tarea.
\end{warningbox}

\subsection{¿Los modelos se comerán a los productos?}

Esta subsección aborda la preocupación principal de las empresas de producto. Guangmi resume la cuestión de fondo como feature system vs learning system. Las empresas de modelos pueden copiar las feature de un producto tradicional; pero si la empresa de producto acumula entorno, datos, flujos de trabajo, estado del usuario, sistema de permisos y ciclo de retroalimentación, una nueva versión del modelo no se la comerá directamente. A la inversa, si el producto solo envuelve al modelo, sin un entorno ni datos propios, el modelo base sí lo absorberá con facilidad.

\begin{importantbox}{Criterio para juzgar la barrera de una aplicación}
Para saber si una aplicación de IA tiene una barrera, no basta con mirar la interfaz y el prompt; hay que ver si posee: un punto de entrada a tareas reales, estado del usuario a largo plazo, retroalimentación verificable, flujos de trabajo propios, control de costos, retorno de datos y costo de sustitución.
\end{importantbox}

\begin{figure}[H]
\centering
\includegraphics[width=0.92\textwidth]{fire-thieves.png}
\caption{Los ladrones del fuego de los modelos: Perplexity, Cursor y Manus llevan las capacidades de los modelos de frontera a flujos de trabajo concretos. Diagrama conceptual propio, elaborado a partir de la conversación de 01:02:45--01:15:11.}
\end{figure}

\begin{knowledgebox}{Lectura de la figura: el valor de los ladrones del fuego está en el contenedor}
Perplexity lleva la capacidad del modelo a la búsqueda y la investigación; Cursor, al entorno de programación; Manus, a la ejecución general. No entrenan ellos mismos el modelo base más fuerte, sino que colocan la capacidad del modelo en un contenedor por el que el usuario está dispuesto a pagar, que usa una y otra vez y que genera retroalimentación.
\end{knowledgebox}

\subsection{Qué debe observar un inversionista}

Esta subsección convierte la cuestión del producto en un juicio de inversión. Guangmi menciona su portafolio ideal y el peso de cada empresa, pero esta nota no lo toma como recomendación de inversión. Lo más importante es el marco de análisis: en las empresas de modelos base se observan el techo de capacidad, la curva de costos y la estabilidad de la organización; en las empresas de aplicaciones, si aprovechan la ventana en que las capacidades del modelo se desbordan hacia otros usos, si acumulan flujos de trabajo y si tienen ingresos reales y costos controlables; en las empresas de infraestructura, si controlan los eslabones clave de entorno, inferencia, entrenamiento, datos y despliegue.

\begin{knowledgebox}{Nota de clase: invertir en tecnología no es apostar por palabras de moda}
Al final, la entrevista enfatiza «crear» en lugar de «moverse en el círculo». En la inversión en IA, esto significa entender las rutas tecnológicas y la cadena de valor: quién eleva la inteligencia, quién la recibe, quién proporciona el entorno, quién controla la retroalimentación y quién solo hace arbitraje de narrativa.
\end{knowledgebox}

\subsection{Resumen de la sección}

La relación entre modelo y producto no se reduce a quién se come a quién. Los modelos fuertes reducen el espacio de los productos delgados, pero también crean oportunidades para nuevos contenedores. La barrera de un producto proviene del entorno, el estado, la retroalimentación, el costo y las tareas del usuario, no de poner una simple ventana de chat encima del modelo.
